% Vietnamese translation for OI Public Library
% Source-File: IOI中国国家候选队论文/2006/唐文斌/唐文斌.doc
\documentclass[11pt]{article}
\usepackage{oipl-vn}
\usepackage{tikz}
\usetikzlibrary{arrows.meta,positioning}

\newcommand{\Oh}{\mathcal{O}}

\title{Bàn về ứng dụng tư tưởng điều chỉnh trong thi tin học}
\author{Tang Wenbin\\Trường Trung học số 1 Thiệu Hưng, Chiết Giang}
\date{2006}

\begin{document}
\maketitle

\origtitle{Bàn về ứng dụng tư tưởng điều chỉnh trong thi tin học}
\sourcepath{IOI China candidate team papers / 2006 / Tang Wenbin.doc}

\renewcommand{\abstractname}{Tóm tắt}
\begin{abstract}
Khi độ khó của bài thi tin học tăng lên và quan hệ dữ liệu trở nên phức tạp,
không phải lúc nào ta cũng tìm được trực tiếp một phương pháp tối ưu. Một cách
tiếp cận tự nhiên là trước hết tạo một phương án bất kỳ hoặc một phương án thô,
sau đó dùng một chuỗi thao tác điều chỉnh, biến đổi và cải tiến để làm phương
án tiến dần tới yêu cầu. Bài viết trình bày tư tưởng ấy qua một số bài toán
tiêu biểu, từ bài dựng nghiệm thỏa điều kiện tới các bài mở không yêu cầu tối
ưu tuyệt đối.
\end{abstract}

\noindent\textbf{Từ khóa:} thi tin học; tư tưởng điều chỉnh; ngẫu nhiên hóa;
cải tiến phương án; bài toán mở.

\section{Dẫn nhập}

``Điều chỉnh'' là một khái niệm rất quen thuộc trong đời sống: nước quá nóng
thì hạ nhiệt, quá lạnh thì tăng nhiệt. Nghĩa chung của nó là thay đổi trạng
thái ban đầu để thích hợp hơn với môi trường và yêu cầu.

Trong khoa học máy tính, tư tưởng này xuất hiện trong nhiều thuật toán. Thuật
toán đơn hình cho quy hoạch tuyến tính liên tục chuyển từ một đỉnh của đa diện
tới đỉnh tốt hơn. Mô phỏng tôi luyện cũng liên tục sinh nghiệm mới, chấp nhận
hoặc loại bỏ nó theo một tiêu chuẩn xác suất. Trong thi tin học, cùng một tinh
thần ấy thường giúp ta xử lý những bài mà lời giải trực tiếp khó tìm: ta tạo
một phương án ban đầu, rồi sửa dần để nó trở thành hợp lệ hoặc tốt hơn.

\section{Điều chỉnh trong bài toán dựng nghiệm}

\subsection{Ví dụ 1: Viễn thông}

Bài \textit{Telecommunication} của Baltic 2001 mô tả hai đảo Bornholm và
Gotland. Mỗi đảo có một số cổng viễn thông; mỗi cổng nối tới một cổng mục tiêu
trên đảo kia. Mỗi cổng được đặt ở một trong hai chế độ: gửi hoặc nhận. Cần
chọn chế độ cho tất cả các cổng sao cho mọi cổng đều hoạt động:
\begin{itemize}
  \item nếu một cổng ở chế độ nhận, phải có ít nhất một cổng ở đảo kia gửi tới
  nó;
  \item nếu một cổng ở chế độ gửi, cổng mục tiêu của nó phải ở chế độ nhận.
\end{itemize}
Giới hạn trong bản gốc là \(1\le n,m\le 50000\).

Một số cổng có thể bị xác định ngay. Chẳng hạn, nếu không có cổng nào gửi tới
một cổng \(A\), thì \(A\) không thể là cổng nhận, nên nó phải là cổng gửi.
Quyết định này lại ép cổng mục tiêu của \(A\) thành cổng nhận. Tuy nhiên,
không phải đồ thị nào cũng có điểm bắt đầu rõ ràng như vậy; có những cấu hình
mọi cổng đều có tiền nhiệm, và ta không thể tiếp tục chỉ bằng suy luận cục bộ.

Cách điều chỉnh đơn giản hơn là bắt đầu bằng một phương án rất thô:
\begin{enumerate}
  \item đặt mọi cổng trên Bornholm ở chế độ gửi;
  \item đặt mọi cổng trên Gotland ở chế độ nhận;
  \item khi còn một cổng nhận vô hiệu trên Gotland, đổi cổng ấy sang chế độ
  gửi và đặt cổng mục tiêu của nó trên Bornholm sang chế độ nhận.
\end{enumerate}

Sau mỗi lần điều chỉnh, số cổng nhận trên Gotland giảm đi một, nên thuật toán
chắc chắn dừng. Khi dừng, mọi cổng nhận trên Gotland đều hợp lệ. Mỗi cổng gửi
trên Gotland đã buộc mục tiêu của nó thành cổng nhận. Mỗi cổng nhận trên
Bornholm cũng nhận được ít nhất một cổng Gotland gửi tới, còn các cổng gửi
trên Bornholm vốn trỏ tới các cổng nhận trên Gotland từ phương án ban đầu.
Vì vậy phương án cuối cùng hợp lệ.

\begin{figure}[ht]
\centering
\begin{tikzpicture}[
  >=Latex,
  port/.style={circle, draw, minimum size=7mm, inner sep=0pt},
  send/.style={port, fill=gray!15},
  recv/.style={port, double},
  node distance=8mm and 38mm,
  every label/.style={font=\scriptsize}
]
\node[send, label=left:{gửi}] (b1) {\(B_1\)};
\node[send, below=of b1, label=left:{gửi}] (b2) {\(B_2\)};
\node[recv, below=of b2, label=left:{nhận}] (b3) {\(B_3\)};

\node[recv, right=of b1, label=right:{nhận}] (g1) {\(G_1\)};
\node[recv, below=of g1, label=right:{nhận}] (g2) {\(G_2\)};
\node[send, below=of g2, label=right:{gửi}] (g3) {\(G_3\)};

\draw[->] (b1) -- (g1);
\draw[->] (b2) -- (g2);
\draw[->] (g3) -- (b3);
\draw[->, dashed] (b3) -- (g2);

\node[above=5mm of b1] {Bornholm};
\node[above=5mm of g1] {Gotland};
\node[below=7mm of b3, align=center, font=\small]
  {Nếu một cổng nhận trên Gotland không có cung gửi tới,\\
   đổi nó sang gửi và ép cổng mục tiêu bên kia thành nhận.};
\end{tikzpicture}
\caption{Một bước điều chỉnh trong bài viễn thông: trạng thái gửi/nhận được
sửa cục bộ nhưng vẫn giữ các ràng buộc đã được kiểm soát.}
\end{figure}

Khi cài đặt, mỗi cổng chỉ cần bị đổi trạng thái nhiều nhất một lần. Ta duy trì
bậc vào của các cổng nhận và dùng hàng đợi cho các cổng nhận vô hiệu. Tổng
thời gian là tuyến tính theo số cổng và số liên kết, đạt đúng bản chất của bài
toán.

\subsection{Ví dụ 2: Chu trình Euler trong đồ thị hỗn hợp}

Cho một đồ thị có cả cạnh có hướng và cạnh vô hướng. Cần định hướng các cạnh
vô hướng sao cho đồ thị thu được có chu trình Euler, hoặc kết luận vô nghiệm.
Không thể thay mỗi cạnh vô hướng bằng hai cung ngược chiều, vì mỗi cạnh chỉ
được đi qua một lần.

Điều kiện cần trước hết là đồ thị nền liên thông. Ngoài ra, tại mỗi đỉnh,
tổng bậc phải chẵn; số cung có hướng đi vào hoặc đi ra cũng không được vượt quá
một nửa tổng bậc cần cân bằng. Nếu các điều kiện này không thỏa, bài toán vô
nghiệm.

Khi các điều kiện cần thỏa, ta có thể định hướng tùy ý tất cả cạnh vô hướng để
tạo một phương án ban đầu. Phương án ấy có thể chưa cân bằng, tức tại một số
đỉnh \(u\) có \(\operatorname{out}(u)\ne\operatorname{in}(u)\). Mục tiêu của
điều chỉnh là làm các đỉnh cân bằng dần.

Nếu có một đỉnh \(a\) với
\[
  \operatorname{out}(a)>\operatorname{in}(a),
\]
ta đi theo các cạnh vô hướng đã được định hướng để tìm một đỉnh \(b\) có
\[
  \operatorname{out}(b)<\operatorname{in}(b).
\]
Sau đó đảo hướng toàn bộ đường đi từ \(a\) tới \(b\). Việc đảo này làm độ lệch
của \(a\) và \(b\) giảm, còn các đỉnh giữa đường không đổi hiệu vào-ra.

Vì tổng bậc vào bằng tổng bậc ra, một vùng chỉ gồm các đỉnh thừa bậc ra không
thể tự khép kín mãi. Do đó từ một đỉnh thừa bậc ra luôn tìm được đường tới một
đỉnh thừa bậc vào. Lặp thao tác đảo đường cho tới khi không còn đỉnh lệch, ta
thu được định hướng Euler.

Ở đây tư tưởng điều chỉnh biến một định hướng bất kỳ thành định hướng hợp lệ.
Ta không cần đoán trực tiếp hướng đúng của từng cạnh; chỉ cần biết phép sửa nào
làm phương án gần điều kiện hơn.

\section{Điều chỉnh trong bài toán mở không tối ưu tuyệt đối}

\subsection{Ví dụ 3: Lắp ráp linh kiện}

Bài này là dạng nộp đáp án. Có \(n\) dây chuyền và \(m\) loại bộ phận. Với mỗi
loại bộ phận có \(n\) linh kiện tương ứng, mỗi linh kiện có thời gian lắp ráp
khác nhau. Cần phân phối các linh kiện vào \(n\) sản phẩm sao cho thời gian
hoàn thành muộn nhất càng nhỏ càng tốt. Trong dữ liệu lớn, \(n\) và \(m\) có
thể tới \(500\).

Có thể xem dữ liệu là một ma trận \(n\times m\). Trong mỗi cột, ta được hoán
vị các phần tử; mục tiêu là làm giá trị lớn nhất trong các tổng hàng càng nhỏ
càng tốt. Bài toán gần với các bài phân hoạch kiểu ba-lô và không phù hợp với
quy hoạch động hay tìm kiếm vét cạn ở kích thước này.

Một tham lam tự nhiên là xử lý từng cột: ở cột hiện tại, gán phần tử lớn nhất
cho hàng có tổng hiện tại nhỏ nhất, phần tử lớn nhì cho hàng có tổng nhỏ nhì,
và tiếp tục như vậy. Cách này hợp lý cục bộ nhưng thường chưa đủ tốt, vì một
quyết định tốt ở cột hiện tại có thể làm mất cân bằng tổng thể.

Tư tưởng điều chỉnh xét trực tiếp một phương án \(A\). Gọi \(S_i\) là tổng
hàng \(i\). Nếu trong cùng một cột có hai phần tử \(A_{r_1,c}\) và \(A_{r_2,c}\)
sao cho sau khi đổi chỗ, hai tổng \(S_{r_1}\) và \(S_{r_2}\) gần nhau hơn, ta
thực hiện đổi chỗ ấy. Khi không còn cặp như vậy, phương án là một cực trị cục
bộ theo phép đổi trong một cột.

Vì cực trị cục bộ phụ thuộc vào phương án ban đầu, ta có thể lặp nhiều lần với
khởi tạo ngẫu nhiên khác nhau, điều chỉnh tới cực trị cục bộ, rồi giữ phương án
tốt nhất. Với bài nộp đáp án, cách này tận dụng tốt việc đề không yêu cầu chứng
minh tối ưu tuyệt đối.

\subsection{Ví dụ 4: Bài toán chia vàng}

Có \(N\le 10000\) hộp vàng, hộp thứ \(i\) có giá trị \(A_i\). Cần chia các hộp
cho \(M\le 1000\) tướng, mỗi hộp không được chia nhỏ, sao cho chênh lệch giữa
người nhận nhiều nhất và ít nhất không vượt quá \(K\). Đề bảo đảm có nghiệm.

Bài toán cũng có bản chất khó khi xét tối ưu hoặc dựng trực tiếp. Một số lời
giải tham lam có thể phát bài từ lớn tới nhỏ cho người đang nhận ít nhất, hoặc
đặt chỉ tiêu riêng khi \(K\le1\), thậm chí dùng một bước ba-lô nhỏ để ghép đủ
chỉ tiêu. Các cách ấy tận dụng cấu trúc dữ liệu, nhưng tương đối rườm rà.

Một thao tác điều chỉnh toàn cục gọn hơn như sau. Với mỗi tướng \(i\), chia
tập vàng hiện có của người đó thành hai phần rời nhau \(S_A(i)\) và \(S_B(i)\).
Có thể chọn một ngưỡng ngẫu nhiên \(C\), rồi đưa dần các hộp vào \(S_A(i)\) cho
tới khi tổng đạt gần \(C\), phần còn lại đưa vào \(S_B(i)\).

Sau khi bỏ thông tin cũ về tướng, xem các phần \(S_A(i)\) và \(S_B(i)\) như
hai dãy độc lập. Ghép phần lớn nhất của dãy trên với phần nhỏ nhất của dãy
dưới, phần lớn nhì với phần nhỏ nhì, và tiếp tục. Theo trực giác sắp xếp đối
ngược, phép ghép này làm các tổng mới cân bằng hơn hoặc ít nhất không xấu đi
rõ rệt. Lặp thao tác với các cách chia ngẫu nhiên thường nhanh chóng tạo được
nghiệm thỏa \(K\).

\begin{figure}[ht]
\centering
\begin{tikzpicture}[
  part/.style={draw, rounded corners=1pt, minimum width=16mm, minimum height=8mm, font=\small},
  >=Latex,
  node distance=10mm
]
\node[part] (a1) {\(S_A(1)\)};
\node[part, right=of a1] (a2) {\(S_A(2)\)};
\node[part, right=of a2] (a3) {\(S_A(3)\)};
\node[left=6mm of a1, font=\small] {sắp giảm};

\node[part, below=16mm of a1] (b1) {\(S_B(1)\)};
\node[part, right=of b1] (b2) {\(S_B(2)\)};
\node[part, right=of b2] (b3) {\(S_B(3)\)};
\node[left=6mm of b1, font=\small] {sắp tăng};

\draw[->] (a1) -- (b3);
\draw[->] (a2) -- (b2);
\draw[->] (a3) -- (b1);

\node[below=8mm of b2, align=center, font=\small]
  {Ghép phần lớn với phần nhỏ, phần lớn nhì với phần nhỏ nhì,\\
   rồi lặp lại trên phương án mới.};
\end{tikzpicture}
\caption{Điều chỉnh toàn cục trong bài chia vàng bằng cách tách mỗi phần thành
\(S_A,S_B\), sắp xếp ngược chiều và ghép lại.}
\end{figure}

Ví dụ này khác ví dụ trước ở quy mô điều chỉnh: thay vì đổi hai phần tử trong
một cột, ta mổ xẻ toàn bộ phương án rồi ghép lại. Cả hai vẫn có cùng tinh thần:
dùng phép biến đổi được thiết kế để kéo phương án về phía cân bằng.

\section{Hai dạng điều chỉnh}

Các ví dụ trên có thể chia thành hai nhóm.

Trong nhóm dựng nghiệm, mục tiêu là tìm một phương án thỏa điều kiện. Ta bắt
đầu từ một phương án chưa chắc hợp lệ rồi điều chỉnh để biến ``không hợp lệ''
thành ``hợp lệ''. Bài viễn thông và bài Euler hỗn hợp thuộc nhóm này.

Trong nhóm bài mở, phương án ban đầu đã hợp lệ hoặc gần hợp lệ, nhưng chất
lượng chưa cao. Ta điều chỉnh để biến ``có nghiệm'' thành ``nghiệm tốt hơn''.
Bài lắp ráp linh kiện và bài chia vàng thuộc nhóm này.

Một phép điều chỉnh hữu ích thường có ba đặc điểm:
\begin{itemize}
  \item có đại lượng đo tiến triển, chẳng hạn số cổng nhận vô hiệu hoặc độ lệch
  vào-ra;
  \item thao tác cục bộ không phá hỏng những phần đã được kiểm soát, hoặc nếu
  có phá thì vẫn làm đại lượng mục tiêu tốt hơn;
  \item có thể cài đặt đủ nhanh bằng cấu trúc dữ liệu phù hợp.
\end{itemize}

\section{Mô phỏng tôi luyện}

Phụ lục của bản gốc nhắc tới mô phỏng tôi luyện như một ví dụ kinh điển của
tư tưởng điều chỉnh. Thuật toán bắt đầu với nhiệt độ \(T\) lớn và một nghiệm
hiện tại \(S\). Ở mỗi vòng, sinh nghiệm mới \(S'\), tính
\[
  \Delta=C(S')-C(S).
\]
Nếu \(\Delta<0\), nghiệm mới tốt hơn và được nhận. Nếu \(\Delta\ge0\), nghiệm
mới vẫn có thể được nhận với xác suất xấp xỉ
\[
  \exp\!\left(-\frac{\Delta}{T}\right).
\]
Khi \(T\) giảm dần, thuật toán chuyển từ khám phá rộng sang cải tiến ổn định.
Khác với điều chỉnh tham lam thuần túy, mô phỏng tôi luyện đôi khi chấp nhận
bước đi xấu hơn để thoát khỏi cực trị cục bộ.

\section{Kết luận}

Tư tưởng điều chỉnh không phải một thuật toán cố định, mà là một cách nhìn.
Khi lời giải trực tiếp khó tìm, ta có thể hỏi:
\begin{itemize}
  \item phương án ban đầu dễ tạo nhất là gì;
  \item điều kiện nào khiến nó sai hoặc chưa tốt;
  \item có phép biến đổi nào làm sai lệch giảm đi mà vẫn dễ kiểm soát không.
\end{itemize}

Trong bài dựng nghiệm, điều chỉnh là quá trình từ không có nghiệm rõ ràng tới
có nghiệm. Trong bài mở, điều chỉnh là quá trình từ có nghiệm tới nghiệm tốt
hơn. ``Từ không đến có, từ có đến tốt'' chính là tinh thần cốt lõi của phương
pháp này.

\section*{Ghi chú về hình và phụ lục}

Phụ lục mã nguồn cho các ví dụ không được chép lại vì đây là các chương trình
mẫu dài, không phải nội dung lý thuyết chính.

\begin{thebibliography}{9}
\bibitem{huang}
Huang Yuanhe, \textit{The Emperor's Riddle (URAL 1144) solution report},
bài tập đội tuyển quốc gia năm 2004.

\bibitem{boi}
Baltic Olympiad in Informatics 2001, problem and solution materials.

\bibitem{algorithm-art}
Liu Rujia và Huang Liang, \textit{Nghệ thuật thuật toán và thi tin học}.
\end{thebibliography}

\end{document}
